\begin{abstract}
소수의 정상 이미지만 사용하는 few-shot 산업 이상 국소화에서는 서포트에 없는 정상 외관 변화도 큰 최근접 거리를 만들 수 있어, 높은 점수만으로 정상 변동과 결함을 구분하기 어렵다.
본 논문은 제한된 서포트에서 비교 단서를 늘리기 위해 동결 백본의 표현 계층과 입력 해상도를 독립 매칭 단위로 확장하고, 그 거리를 정상 서포트 기준으로 통합하는 \textbf{\method}를 제안한다.
\method는 각 관측의 거리를 정상 참조 거리의 상위 분위수(Q95)로 나누어 단위를 맞추고, 깊은 계층의 평균에 얕은 계층을 작은 가중치로 더한다.
해상도 축에서는 위치별 점수 순위에 따라 Rényi 가중치를 적용하여 산술 평균에서 생기는 국소 반응의 희석을 줄인다.
백본의 추가 학습이나 쿼리 시점 파라미터 갱신은 사용하지 않는다.
MVTec AD 2에서 \method는 1--16-shot의 모든 조건에서 비교 방법보다 높은 \AUPRO, \pAUROC, \IAUROC를 기록했으며, 1-shot \AUPRO는 52.18\,\%로 차순위보다 5.87\,pp 높다.
구성 요소 분석에서는 다해상도 관측과 관측별 독립 매칭이 가장 큰 향상을 주었고, Q95 보정 아래에서 Rényi 집계가 추가 향상을 주었다.
MVTec AD와 VisA에서도 비교 방법과 대등한 성능을 보였다.
\end{abstract}
